% GENERATED FILE. Edit canonical lessons, then run npm run generate:books.
% canonical-source-hash: fnv1a64:619e7b61dda466e2
% canonical-lessons: PT-C439-nordeste, PT-C439-noroeste, PT-C439-sudeste, PT-C439-sudoeste, PT-C439-marco

\chapter{Norte, sul, leste e oeste}
\label{ch:pt-norte-sul-leste-e-oeste}
\hlchaptermodality{439}

\begin{chapteropening}
\textbf{By the end of this chapter:} \emph{I can place things on the compass: north-east, north-west, south-east, south-west and a landmark.}

Complete the last lesson of chapter 439: I can place things on the compass: north-east, north-west, south-east, south-west and a landmark.
\end{chapteropening}

\section[o nordeste]{o nordeste — the north-east}

\label{lesson:PT-C439-nordeste}



\begin{quote}
Before the new one: say the Portuguese for to run alongside, then the Portuguese for to come out.
\end{quote}

\subsection*{You'll want to know: o nordeste}
\textbf{o nordeste} — \textquotedblleft{}the north-east\textquotedblright{}.

\begin{quote}
A aldeia fica no nordeste do país.
\end{quote}

\begin{quote}
\emph{The village is in the north-east of the country.}
\end{quote}

Written with a capital as a region's name.

\begin{cousinweb}
\emph{nordeste} came from French \emph{nord-est} (EtymDB), from \emph{nord} 'north' and \emph{est} 'east', words of Germanic origin.
\end{cousinweb}

\begin{grammarlens}[title={What you've built}]
One more word for telling, explaining and describing at B1.
\end{grammarlens}

\subsection*{Your turn}
\begin{itemize}
\raggedright
  \item \emph{Say it:} \emph{o nordeste}
  \item \emph{Say it:} \emph{o nordeste}, once more
  \item \emph{Say it:} a sentence of your own with \emph{o nordeste}
  \item \emph{Recall:} say the Portuguese for to run alongside, then the Portuguese for to come out, then say \emph{o nordeste} again
\end{itemize}

\begin{tcolorbox}[breakable,colback=teal!4,colframe=teal!35!black,title={Before you move on}]
What does o nordeste mean? (\textquotedblleft{}The north-east\textquotedblright{}.) Say it once more.
\end{tcolorbox}
\section[o noroeste]{o noroeste — the north-west}

\label{lesson:PT-C439-noroeste}



\begin{quote}
Before the new one: say the Portuguese for to come out, then the Portuguese for the north-east.
\end{quote}

\subsection*{You'll want to know: o noroeste}
\textbf{o noroeste} — \textquotedblleft{}the north-west\textquotedblright{}.

\begin{quote}
O vento vem de noroeste.
\end{quote}

\begin{quote}
\emph{The wind is coming from the north-west.}
\end{quote}

\emph{Vento de noroeste} — a north-westerly.

\begin{cousinweb}
\emph{noroeste} is built from \emph{norte} and \emph{oeste} on the model of French \emph{nord-ouest}; like \emph{nordeste} (EtymDB), it rests on Germanic words for 'north' and 'west'.
\end{cousinweb}

\begin{grammarlens}[title={What you've built}]
One more word for telling, explaining and describing at B1.
\end{grammarlens}

\subsection*{Your turn}
\begin{itemize}
\raggedright
  \item \emph{Say it:} \emph{o noroeste}
  \item \emph{Say it:} \emph{o noroeste}, once more
  \item \emph{Say it:} a sentence of your own with \emph{o noroeste}
  \item \emph{Recall:} say the Portuguese for to come out, then the Portuguese for the north-east, then say \emph{o noroeste} again
\end{itemize}

\begin{tcolorbox}[breakable,colback=teal!4,colframe=teal!35!black,title={Before you move on}]
What does o noroeste mean? (\textquotedblleft{}The north-west\textquotedblright{}.) Say it once more.
\end{tcolorbox}
\section[o sudeste]{o sudeste — the south-east}

\label{lesson:PT-C439-sudeste}



\begin{quote}
Before the new one: say the Portuguese for the north-east, then the Portuguese for the north-west.
\end{quote}

\subsection*{You'll want to know: o sudeste}
\textbf{o sudeste} — \textquotedblleft{}the south-east\textquotedblright{}.

\begin{quote}
A tempestade vai para sudeste.
\end{quote}

\begin{quote}
\emph{The storm is moving south-east.}
\end{quote}

Also spelled \emph{sueste} in old texts.

\begin{cousinweb}
\emph{sudeste} came from French \emph{sud-est} (EtymDB), from \emph{sud} 'south' and \emph{est} 'east', both of Germanic origin.
\end{cousinweb}

\begin{grammarlens}[title={What you've built}]
One more word for telling, explaining and describing at B1.
\end{grammarlens}

\subsection*{Your turn}
\begin{itemize}
\raggedright
  \item \emph{Say it:} \emph{o sudeste}
  \item \emph{Say it:} \emph{o sudeste}, once more
  \item \emph{Say it:} a sentence of your own with \emph{o sudeste}
  \item \emph{Recall:} say the Portuguese for the north-east, then the Portuguese for the north-west, then say \emph{o sudeste} again
\end{itemize}

\begin{tcolorbox}[breakable,colback=teal!4,colframe=teal!35!black,title={Before you move on}]
What does o sudeste mean? (\textquotedblleft{}The south-east\textquotedblright{}.) Say it once more.
\end{tcolorbox}
\section[o sudoeste]{o sudoeste — the south-west}

\label{lesson:PT-C439-sudoeste}



\begin{quote}
Before the new one: say the Portuguese for the north-west, then the Portuguese for the south-east.
\end{quote}

\subsection*{You'll want to know: o sudoeste}
\textbf{o sudoeste} — \textquotedblleft{}the south-west\textquotedblright{}.

\begin{quote}
A praia fica a sudoeste da cidade.
\end{quote}

\begin{quote}
\emph{The beach is south-west of the town.}
\end{quote}

\emph{A costa sudoeste} — the south-west coast.

\begin{cousinweb}
\emph{sudoeste} came from French \emph{sud-ouest} 'south-west' (EtymDB), built on Germanic words for 'south' and 'west'.
\end{cousinweb}

\begin{grammarlens}[title={What you've built}]
One more word for telling, explaining and describing at B1.
\end{grammarlens}

\subsection*{Your turn}
\begin{itemize}
\raggedright
  \item \emph{Say it:} \emph{o sudoeste}
  \item \emph{Say it:} \emph{o sudoeste}, once more
  \item \emph{Say it:} a sentence of your own with \emph{o sudoeste}
  \item \emph{Recall:} say the Portuguese for the north-west, then the Portuguese for the south-east, then say \emph{o sudoeste} again
\end{itemize}

\begin{tcolorbox}[breakable,colback=teal!4,colframe=teal!35!black,title={Before you move on}]
What does o sudoeste mean? (\textquotedblleft{}The south-west\textquotedblright{}.) Say it once more.
\end{tcolorbox}
\section[o marco]{o marco — the landmark; the milestone; the boundary stone}

\label{lesson:PT-C439-marco}



\begin{quote}
Before the new one: say the Portuguese for the south-east, then the Portuguese for the south-west.
\end{quote}

\subsection*{You'll want to know: o marco}
\textbf{o marco} — \textquotedblleft{}the landmark; the milestone; the boundary stone\textquotedblright{}.

\begin{quote}
A torre é um marco da cidade.
\end{quote}

\begin{quote}
\emph{The tower is a landmark of the town.}
\end{quote}

\emph{Um marco do correio} — a pillar box. \emph{Um marco histórico} — a milestone.

\begin{cousinweb}
\emph{marco} comes from Medieval Latin \emph{marcus}, from Germanic \emph{markō} 'boundary, border' (EtymDB): first a stone that marked the edge of a field.
\end{cousinweb}

\begin{grammarlens}[title={What you've built}]
One more word for telling, explaining and describing at B1.
\end{grammarlens}

\subsection*{Your turn}
\begin{itemize}
\raggedright
  \item \emph{Say it:} \emph{o marco}
  \item \emph{Say it:} \emph{o marco}, once more
  \item \emph{Say it:} a sentence of your own with \emph{o marco}
  \item \emph{Recall:} say the Portuguese for the south-east, then the Portuguese for the south-west, then say \emph{o marco} again
\end{itemize}

\begin{tcolorbox}[breakable,colback=teal!4,colframe=teal!35!black,title={Before you move on}]
What does o marco mean? (\textquotedblleft{}The landmark; the milestone; the boundary stone\textquotedblright{}.) Say it once more.
\end{tcolorbox}
